%\vspace{-0.3cm}
\section{Related Work}
%\vspace{-0.4cm}
% In this work, we borrowed concepts from Natural Language Processing (NLP), fairness, and interpretability. Thus, we discuss various insights from each of these concepts in this section.
% \subsection{Fairness}
\para{Fairness.} The research in fairness concerns itself with various topics, such as defining fairness metrics, proposing solutions for bias mitigation, and analyzing existing harms in various systems~\cite{10.1145/3457607}. In this work, we utilized different metrics that were introduced previously, such as statistical parity~\cite{Dwork:2012:FTA:2090236.2090255}, equality of opportunity and equalized odds~\cite{NIPS2016_9d268236}, to measure the amount of bias. We also used different bias mitigation strategies to compare against our mitigation strategy, such as FCRL~\cite{gupta2021controllable}, CVIB~\cite{NEURIPS2018_415185ea}, MIFR~\cite{pmlr-v89-song19a}, adversarial forgetting~\cite{advforg}, MaxEnt-ARL~\cite{roy2019mitigating}, and LAFTR~\cite{pmlr-v80-madras18a}.  We also utilized concepts and datasets that were analyzing existing biases in NLP systems, such as~\cite{de2019bias} which studied the existing biases in NLP systems on the occupation classification task on the bios dataset.
% \subsection{Interpretability}

\para{Interpretability.} In this work, we introduced an attribution framework based on the attention weights that can analyze fairness and accuracy of the models at the same time and reason about the importance of each feature on fairness and accuracy. There is a body of work in NLP literature that tried to analyze the effect of the attention weights on interpretability of the model~\cite{wiegreffe-pinter-2019-attention,jain-wallace-2019-attention,serrano-smith-2019-attention}. Other work also utilized attention weights to define an attribution score to be able to reason about how transformer models such as BERT work~\cite{hao2021self}. Notice that although \citet{jain-wallace-2019-attention} claim that attention might not be explanation, a body of work has proved otherwise including~\cite{wiegreffe-pinter-2019-attention} in which authors directly target the work in \citet{jain-wallace-2019-attention} and analyze in detail the problems associated with this study. In our work, we also find that attention can be useful and can extract meaningful information which can be beneficial in many aspects. In addition,~\citet{NEURIPS2020_92650b2e} analyze the effect of the attention weights in transformer models for bias analysis in language models. However, their approach is different and has a more causal take on investigating the bias. Their study is specific to language models and does not necessarily apply to broader tasks and existing fairness definitions. Aside from interpretability and fairness, we utilized concepts from the NLP literature for designing our attention-based model that can be applicable to tabular data~\cite{NIPS2017_3f5ee243,zhou-etal-2016-attention}.
 